\subsection{Cúbits}
Como introduce \textcite{diracBraket}, los cúbits se escriben normalmente en la llamada notación de Dirac o notación bra-ket.
Cada cúbit se representa como un ket, de la forma $\ket{a}$, que denota un vector en $\Q$.

Por ejemplo, los estados de la base canónica se escriben como $\ket{0}$ y $\ket{1}$, llamados ket 0 y ket 1, respectivamente (véase~\cite{NielsenChuang}).
Así, cualquier cúbit puede describirse como una combinación lineal de estos dos estados de la base: $\ket{a} =a_0\ket{0} +a_1 \ket{1}$, donde $a_0, a_1 \in \mathbb{C}$ con $|a_0|^{2}+|a_1|^{2}=1$.

\begin{notation}
  Para cualquier cúbit $\ket{a}$, se usa la siguiente notación:
  \[a_{0}=\braket{0}{a}\text{ y }a_{1}=\braket{1}{a},\]
\end{notation}
por lo que el cúbit $\ket{a}$ puede escribirse como:
\begin{equation}
  \label{eq:cúbit-coordenadas}
  \ket{a}=a_0\ket{0}+a_1\ket{1}\,.
\end{equation}

Para otras bases, como por ejemplo $\set{\ket{+}, \ket{-}}$, donde $\ket{+}=\frac{\ket{0}+\ket{1}}{\sqrt{2}}$ y $\ket{-}=\frac{\ket{0}-\ket{1}}{\sqrt{2}}$, un cúbit $\ket{a}$ puede describirse como:
\begin{equation}
  \label{eq:cúbit-coordenadas-h}
  \ket{a}=a_+\ket{+}+a_-\ket{-}\,.
\end{equation}

\subsection{Kets}

En general, para $n$ cúbits $\ket{a_1},\dots,\ket{a_n}$, el ket compuesto $\ket{a}=\ket{a_1\cdots a_n}\in\Q_n$ puede escribirse como se muestra a continuación:
\begin{equation}
  \label{eq:ncubit-coordenadas}
  \ket{a}=a_{k0}\ket{a_1\cdots a_{k-1}0a_{k+1}\cdots a_n}+a_{k1}\ket{a_1\cdots a_{k-1}1a_{k+1}\cdots a_n}.
\end{equation}

Por ejemplo, para dos cúbits $\ket{a}$ y $\ket{b}$, las siguientes expresiones son iguales:
\begin{align*}
  \ket{ab} & =a_0\ket{0b}+a_1\ket{1b}                                            \\
           & =b_0\ket{a0}+b_1\ket{a1}                                            \\
           & =a_0 b_0\ket{00}+a_0 b_1\ket{01}+a_1 b_0\ket{10}+a_1 b_1\ket{11}\,.
\end{align*}

\begin{notation}
  Para cualquier ket $\ket{a}=\ket{a_1\dots a_n}$ de $n$ cúbits, dado un cúbit $\ket{b}$, se usa la siguiente notación:
  \begin{equation}
    \label{eq:ncubit-only-one-state}
    \stateket{a}{k}{b} = \ket{a_1\cdots a_{k-1}b a_{k+1}\cdots a_n},
  \end{equation}
  donde $k \in \mathbb{N}_n$.
\end{notation}

En consecuencia, la expresión~(\ref{eq:ncubit-coordenadas}) puede escribirse como sigue:
\begin{equation}
  \label{eq:canonical-state}
  \ket{a}= a_{k0}\stateket{a}{k}{0}+a_{k1}\stateket{a}{k}{1}.
\end{equation}

Por ejemplo, si $\ket{a_1},\ket{a_2}$ y $\ket{a_3}$ son tres cúbits, y $\ket{a} = \ket{a_1 a_2 a_3}$, entonces:
\begin{itemize}
  \item $\stateket{a_1}{1}{0}=\ket{0}$ y $\stateket{a_1}{1}{1}=\ket{1}$.
  \item $\stateket{a_1 a_2}{2}{0}=\ket{a_1 0}$.
  \item $\ket{a}=a_{20}\ket{a_1 0 a_3}+a_{21}\ket{a_1 1 a_3}=a_{20}\stateket{a}{2}{0}+a_{21}\stateket{a}{2}{1}$.
  \item $\stateket{00a_3}{3}{1}=\ket{001}$.
\end{itemize}

\begin{notation}
  Para una notación más general basado en conjuntos de índices, dado el conjunto $\Gamma=\set{k_1,\dots,k_l}$ y $\gamma\in\B^l$ denotamos por:
  \[
    \stateket{a}{\Gamma}{\gamma}=((1-\delta_{1\Gamma})\ket{a_1}+\delta_{1\Gamma}\ket{\ext{\gamma}{\Gamma}_1})\otimes\dots\otimes ((1-\delta_{n\Gamma})\ket{a_n}+\delta_{n\Gamma}\ket{\ext{\gamma}{\Gamma}_n})\,.
  \]
\end{notation}

Veamos exactamente que valor toma $(1-\delta_{k\Gamma})\ket{a_k}+\delta_{k\Gamma}\ket{\ext{\gamma}{\Gamma}_k}$:
\begin{itemize}
  \item Si $k\in\Gamma\so\delta_{k\Gamma}=1$ y $\ext{\gamma}{\Gamma}_k=\gamma_k$, por lo tanto
        \begin{equation*}
          (1-\delta_{k\Gamma})\ket{a_k}+\delta_{k\Gamma}\ket{\ext{\gamma}{\Gamma}_k}=(1-1)\ket{a_k}+\ket{\gamma_k}=\ket{\gamma_k}\,.
        \end{equation*}
  \item Si $k\notin\Gamma\so\delta_{k\Gamma}=0$ y $\ext{\gamma}{\Gamma}_k=0$, por lo tanto
        \begin{equation*}
          (1-\delta_{k\Gamma})\ket{a_k}+\delta_{k\Gamma}\ket{\ext{\gamma}{\Gamma}_k}=(1-0)\ket{a_k}+0\ket{0}=\ket{a_k}\,.
        \end{equation*}
\end{itemize}

Por ejemplo, si $\Gamma=\set{1,3,7}\subseteq \N_8$, $\gamma=(1,0,1)\in\B^3$ y $\ket{a}=\ket{0}\in\Q_8$ entonces
\[
  \ket{a}^{\Gamma}_{\gamma}=\ket{10000010}.
\]

La notación anterior sustituye los cúbits de las posiciones de $\Gamma$ por estados de la base computacional. Del mismo modo se pueden sustituir por estados arbitrarios de un cúbit, como $\ket{+}$ o $\ket{-}$.

\begin{notation}
  \label{not:ket-estados}
  Sean $k_1,\dots,k_l\in\N_n$ índices distintos y $\ket{b_1},\dots,\ket{b_l}\in\Q$ cúbits. Se denota por $\stateket{a}{k_1\cdots k_l}{b_1\cdots b_l}$ el ket que resulta de sustituir en $\ket{a}=\ket{a_1\cdots a_n}$ el cúbit $\ket{a_{k_m}}$ por $\ket{b_m}$, para $m=1,\dots,l$, y dejar los demás igual. Si $\Gamma=\set{k_1,\dots,k_l}$ con $k_1<\dots<k_l$, también se escribe $\stateket{a}{\Gamma}{b_1\cdots b_l}$.
\end{notation}

Cuando $\ket{b_m}=\ket{\gamma_m}$ con $\gamma\in\B^l$ se recupera la notación anterior, y cuando $\Gamma=\set{k}$ se recupera~\eqref{eq:ncubit-only-one-state}. Por ejemplo, si $\ket{a}=\ket{a_1a_2a_3}$, entonces $\stateket{a}{13}{+-}=\ket{+\,a_2\,-}$ y $\stateket{a}{31}{+-}=\ket{-\,a_2\,+}$.

Aplicando~\eqref{eq:canonical-state} a cada índice, un ket producto se descompone según los valores de los cúbits $k_1,\dots,k_l$:
\begin{equation}
  \label{eq:descomposicion-indices}
  \ket{a}=\sum_{\gamma\in\B^l} a_{k_1\gamma_1}\cdots a_{k_l\gamma_l}\,\stateket{a}{k_1\cdots k_l}{\gamma_1\cdots\gamma_l}.
\end{equation}

